\section*{Chapitre \ref{ch:m3:maximal-extractable-value} --- Valeur extractible maximale}

\begin{solution}{exo:m3:maximal-extractable-value:1}
Tout \omterm{def:m3:maximal-extractable-value:front}{front-running} fait tomber le montant reçu par la victime sous la cotation, si bien que la transaction échouerait et que l'attaquant ne conserverait que
sa position de \omterm{def:m3:maximal-extractable-value:front}{front-running}. Les portefeuilles tolèrent un certain écart, car le pool peut évoluer pour des raisons innocentes entre la soumission et
l'inclusion (autres transactions, arbitrage), et une transaction sans tolérance échouerait souvent.
\end{solution}

\begin{solution}{exo:m3:maximal-extractable-value:2}
La liquidation qui suit une mise à jour de l'\omterm{def:m3:blockchains-for-traders:contract}{oracle} est un \omterm{def:m3:maximal-extractable-value:front}{back-running}~; l'arbitrage qui suit un gros échange est un \omterm{def:m3:maximal-extractable-value:front}{back-running}~; l'achat qui
précède un gros achat est un \omterm{def:m3:maximal-extractable-value:front}{front-running}.
\end{solution}

\begin{solution}{exo:m3:maximal-extractable-value:3}
Pour le constructeur~: le proposeur ne peut ni voir ni modifier les transactions du bloc avant de s'engager à le proposer. Pour le
proposeur~: l'offre qu'il accepte est valide et sera payée s'il propose le bloc.
\end{solution}

\begin{solution}{exo:m3:maximal-extractable-value:4}
Les deux doublent~: de 5\,067~USD à 10\,130~USD de gain brut, et de 1.56 à 3.12 ether de perte.
\end{solution}

\begin{solution}{exo:m3:maximal-extractable-value:5}
$10\,129.57 \times 10\% - 20 = $ 992.96~USD.
\end{solution}

\begin{solution}{exo:m3:maximal-extractable-value:6}
Les canaux envoient leurs flux aux constructeurs auxquels ils font confiance ou qui les rémunèrent~; les constructeurs disposant du plus grand \omterm{def:m3:maximal-extractable-value:private}{flux d'ordres privé} peuvent construire les blocs les plus
rentables, remporter davantage d'enchères et attirer davantage de flux, dans une boucle de rétroaction.
\end{solution}

\begin{solution}{exo:m3:maximal-extractable-value:7}
Une hausse d'environ 0.54~\% (jusqu'à 3\,016.20)~: l'arbitrage achète pour environ 17\,900~USD d'ether, pour un profit d'environ 21~USD. En dessous des
frais du pool de 0.30~\%, il n'existe aucun arbitrage.
\end{solution}

\begin{solution}{exo:m3:maximal-extractable-value:8}
Le canal empêche les \omterm{def:m3:maximal-extractable-value:sandwich}{attaques sandwich}, mais l'échange déplace toujours le pool, dont la valeur de \omterm{def:m3:maximal-extractable-value:front}{back-running} est partagée ou conservée par le canal et
les constructeurs~; les \omterm{def:m3:blockchains-for-traders:gas}{frais de priorité} et les conditions propres au canal sont des coûts~; et le trader fait confiance au canal et à ses constructeurs pour ne pas utiliser
ces flux. L'exposition au MEV est réduite et déplacée, non supprimée.
\end{solution}

\begin{solution}{pb:m3:maximal-extractable-value:1}
\textbf{1.} 311.62 ether, pour un prix moyen de 3\,209.03 dollars~: environ 7~\% au-dessus du prix du pool, dont 0.30~\% de frais et
le reste étant l'impact propre de l'échange. \textbf{2.} 310.06 et 308.50 ether. \textbf{3.} 38\,912~USD et 78\,119~USD. \textbf{4.} Son
gain croît avec le \omterm{def:m3:maximal-extractable-value:front}{front-running} jusqu'au point où la transaction de la victime échouerait. \textbf{5.} Rien~: aucun \omterm{def:m3:maximal-extractable-value:front}{front-running} n'est réalisable.
\textbf{6.} 5\,067~USD et 10\,130~USD. \textbf{7.} 1.56 et 3.12 ether, soit environ 4\,674~USD et 9\,349~USD à 3\,000. \textbf{8.}
5\,046.72~USD et 10\,109.57~USD sans paiement~; 486.67~USD et 992.96~USD en versant 90~\%. \textbf{9.} Le constructeur, et à travers son
offre le \omterm{def:m3:blockchains-for-traders:pow}{validateur} proposant. \textbf{10.} Chaque \omterm{def:m3:maximal-extractable-value:mev}{searcher} qui voit la même opportunité y enchérit~; la meilleure offre approche
la totalité de la valeur. \textbf{11.} Juste au-dessus du glissement attendu de l'échange, dû aux autres transactions avant l'inclusion, de quelques points de base, ou
l'envoyer en privé~; 0.5~\% fait un cadeau pouvant aller jusqu'à 4\,700~USD. \textbf{12.} Chaque morceau déplace moins le pool, si bien que chaque tolérance peut être
plus serrée et chaque \omterm{def:m3:maximal-extractable-value:sandwich}{attaque sandwich} plus petite~; le trader supporte le risque de prix sur plusieurs blocs et dix fois plus de \omterm{def:m3:blockchains-for-traders:gas}{gas}. \textbf{13.} Elle masque
l'échange aux \omterm{def:m3:maximal-extractable-value:mev}{searchers} publics et peut rembourser la valeur de \omterm{def:m3:maximal-extractable-value:front}{back-running}~; elle ne supprime ni l'impact sur le prix ni la confiance accordée au canal.
\textbf{14.} Les solveurs ou fillers se disputent l'exécution à un prix unique, et au sein d'un lot l'ordre n'a pas d'importance~; le pire
prix du trader reste sa limite signée. \textbf{15.} Repérer dans le bloc la même adresse achetant le même \omterm{def:m3:blockchains-for-traders:key}{jeton} juste avant et
le revendant juste après, dans le même pool (\texttt{find\_sandwiches}), et comparer l'exécution à la cotation. \textbf{16.} Elle exploite
la connaissance d'un ordre en attente pour passer avant lui, ce qui constitue le cœur du \omterm{def:m3:maximal-extractable-value:front}{front-running}, mais l'attaquant ne doit à la victime aucun
devoir, comme le ferait un courtier~; la question de son illégalité dépend de la juridiction et des faits, et elle est contestée. \textbf{17.} Ils
s'appuient sur des informations rendues publiques après la transaction et n'aggravent pas le prix de l'utilisateur~; ils corrigent les prix. \textbf{18.} Que la
séparation des rôles repose sur des logiciels pouvant comporter des failles, et que les participants qui contrôlent des \omterm{def:m3:blockchains-for-traders:pow}{validateurs} peuvent les exploiter~;
la confiance s'est déplacée vers les relais et leur code. \textbf{19.} \emph{Résultat nommé~:} 5\,046.72~USD nets et 1.56 ether perdus pour une tolérance de 0.5~\%~;
10\,109.57~USD nets et 3.12 ether perdus pour une tolérance de 1~\%. \textbf{20.} Une option écrite sur le prix de l'échange, exerçable par quiconque ordonne
le bloc.
\end{solution}

\begin{solution}{iq:m3:maximal-extractable-value:1}
Un achat avant et une revente après l'échange de la victime dans le même pool, bornés par la \omterm{def:m3:maximal-extractable-value:sandwich}{tolérance de glissement} de la victime. Se protéger par des
tolérances serrées, des canaux privés, le fractionnement des grands échanges et des mécanismes par intentions ou par lots.
\iqlookfor{la tolérance comme borne, et le \omterm{def:m3:maximal-extractable-value:private}{flux d'ordres privé}.}
\end{solution}

\begin{solution}{iq:m3:maximal-extractable-value:2}
Les \omterm{def:m3:maximal-extractable-value:pbs}{constructeurs de blocs} assemblent les blocs et enchérissent pour l'inclusion~; les proposeurs retiennent la meilleure offre sans voir le contenu~; les relais conservent le
contenu des blocs jusqu'à l'engagement du proposeur (un commit-and-reveal) et vérifient les offres, afin qu'aucune des parties ne puisse tromper l'autre.
\iqlookfor{le commit and reveal et qui fait confiance à qui.}
\end{solution}

\begin{solution}{iq:m3:maximal-extractable-value:3}
Le montant reçu par la victime diminue avec le \omterm{def:m3:maximal-extractable-value:front}{front-running}~; posons que le montant obtenu après un \omterm{def:m3:maximal-extractable-value:front}{front-running} $a$ est égal à $(1 - \tau)$ fois la cotation, et
résolvons pour $a$ (sous forme fermée pour le produit constant, ou par recherche dichotomique en arithmétique entière).
\iqlookfor{monotonie et contrainte active.}
\end{solution}

\begin{solution}{iq:m3:maximal-extractable-value:4}
C'est le même flux vu sous deux angles~: le profit de l'arbitragiste qui ramène le pool obsolète au prix externe est la
perte du pool due au rééquilibrage, nette des frais.
\iqlookfor{le gain de l'arbitragiste égale la perte du pool.}
\end{solution}

\begin{solution}{iq:m3:maximal-extractable-value:5}
Sur le plan opérationnel~: bugs dans les bundles qui font perdre de l'argent, gestion des clés, dépendance aux relais et aux constructeurs, transactions annulées et \omterm{def:m3:blockchains-for-traders:gas}{gas}.
Sur le plan juridique~: les stratégies qui nuisent aux utilisateurs (\omterm{def:m3:maximal-extractable-value:sandwich}{attaques sandwich}) peuvent être traitées comme une fraude ou une manipulation de marché~; les \omterm{def:m3:blockchains-for-traders:pow}{validateurs} ou les exploits
d'infrastructure exposent à un risque pénal, comme le montre l'affaire des relais de 2023.
\iqlookfor{distinction entre le \omterm{def:m3:maximal-extractable-value:front}{back-running} bénin et les stratégies nuisibles.}
\end{solution}

\begin{solution}{iq:m3:maximal-extractable-value:6}
Les utilisateurs envoient leurs transactions à un nœud de confiance qui ne révèle que des indications (pool, sens, mais ni taille ni limite)~; les \omterm{def:m3:maximal-extractable-value:mev}{searchers} enchérissent pour le
droit de faire du \omterm{def:m3:maximal-extractable-value:front}{back-running} avec des bundles que le nœud simule~; le nœud transmet le meilleur bundle à condition que l'utilisateur perçoive une
part convenue, et seuls les back-runs sont acceptés. Renforcer le dispositif par des offres scellées, plusieurs constructeurs et des audits.
\iqlookfor{indications et non ordres, back-runs seuls, et conditions de paiement.}
\end{solution}
